% ============================================================================
% CAPABILITY / ROUTING analysis — STAGING draft (5-block layered structure).
% Anchors the paper's claim on the CONSERVATIVE reading; reports the
% pre-registered result verbatim, then transparently supersedes it.
% Numbers from the full 1300-image RESISC45 re-benchmark (2B=0.394, 3B=0.456,
% 7B=0.549; per-class n=19-41, mean 28.9).
% ============================================================================

% ---- Block 1: pre-registered test + literal result ----
We ask whether any scene categories require the split-only 7B configuration,
that is, whether a single-satellite model is insufficient for them. To avoid
selecting categories after seeing the data, we fixed the test in advance: a
category requires 7B if the 3B model scores at most 0.35, the 7B model scores
at least 0.70, and the difference is significant under a two-sided Fisher exact
test after Benjamini-Hochberg correction across all 45 categories. On the full
1300-image benchmark (19 to 41 images per category, mean 28.9), five categories
pass: parking\_lot, freeway, mobile\_home\_park, industrial\_area, and mountain.

% ---- Block 2: non-monotonicity finding ----
Four of these five expose a weakness of the test rather than a capability gap.
On parking\_lot, freeway, mobile\_home\_park, and industrial\_area the accuracy
is not monotone in model size: the smaller 2B model scores 0.79 to 0.89, the 3B
model drops to 0.04 to 0.34, and 7B recovers to 0.78 to 0.83. This is a
category-specific weakness of the 3B model, not a sign that the category needs
7B. Because 2B also runs on a single satellite and costs less, the best
single-satellite option already serves these four categories.

% ---- Block 3: max(2B,3B) re-analysis, LABELED as a pre-registration deviation ----
The pre-registered test used 3B as the single-satellite baseline. The quantity
relevant to the scheduler is instead the best accuracy reachable by any
single-satellite model, $\max(\text{2B},\text{3B})$. We therefore report a
stricter reanalysis, which we mark as a deviation from the pre-registration: a
category requires 7B only if both single-satellite models fail
($\max(\text{2B},\text{3B})\le 0.35$) while 7B reaches at least 0.70. Under this
criterion only mountain qualifies (2B 0.19, 3B 0.33, 7B 0.78), and it passes
Benjamini-Hochberg but not the stricter Bonferroni threshold. The aggregate
accuracy gain of 7B over 3B is 0.093.

% ---- Block 4: memory-necessity (main pillar) + routing pivot ----
Two conclusions follow. First, the split pipeline is required for a reason that
does not depend on any per-category result: a 4-bit 7B model with its vision
encoder and key-value cache does not fit in the memory left on an 8\,GB
satellite after the sensing pipeline, so the two-satellite split is the only way
to run 7B at all. Second, the accuracy data show that model size alone is a poor
guide to which configuration to use. Most categories are served well by a
single-satellite model, one category needs 7B, and several categories on which
3B fails are recovered by the smaller 2B. Selecting a configuration is therefore
not "larger is better" but a routing problem over heterogeneous categories: the
scheduler reasons over the measured per-configuration quality of each source
rather than over model size, and it is this routing, together with energy
sustainability, that the proposed method provides.

% ---- Block 5: statistical caveats ----
Two caveats bound these claims. The per-category sample is small (19 to 41
images), so a single per-category accuracy carries a 95\% confidence interval of
about $\pm0.18$; per-category necessity claims are fragile, and we rest the case
for the split on the memory constraint rather than on any single category. The
aggregate gain of 0.093 over 1300 paired images is likely significant, but
significance is not necessity: it shows that 7B is better on average, not that
any category cannot be served by a single satellite. We treat mountain as the
one suggestive case and do not generalize from it.
